% ============================================================
\chapter{Transactions et Contr\^ole de Concurrence}
\label{ch:transactions}
% ============================================================

Imaginez un virement bancaire~: on d\'ebite un compte et on cr\'edite un autre. Que se passe-t-il si le syst\`eme tombe en panne entre les deux op\'erations~? L'argent a disparu. C'est pour \'eviter ce genre de catastrophe que Jim Gray, dans les ann\'ees 1970 chez IBM, a formalis\'e la notion de \emph{transaction} et les propri\'et\'es ACID (Atomicit\'e, Coh\'erence, Isolation, Durabilit\'e). Une transaction est un ``tout ou rien'' logique~: soit toutes ses op\'erations r\'eussissent, soit aucune n'a d'effet. Quand plusieurs transactions s'ex\'ecutent simultan\'ement, le contr\^ole de concurrence (verrouillage, estampillage, MVCC) garantit que le r\'esultat est \'equivalent \`a une ex\'ecution s\'equentielle. Gray recevra le prix Turing en 1998 pour ces travaux fondamentaux.

\begin{intuition}[title=\textbf{Id\'ee centrale}]
Une transaction est une unit\'e logique de travail qui doit \^etre ex\'ecut\'ee
de mani\`ere atomique, coh\'erente, isol\'ee et durable (ACID). Le contr\^ole
de concurrence garantit que l'ex\'ecution simultan\'ee de plusieurs
transactions produit un r\'esultat correct, comme si elles avaient \'et\'e
ex\'ecut\'ees s\'equentiellement.
\end{intuition}

% ------------------------------------------------------------
\section{Notion de transaction}
% ------------------------------------------------------------

\begin{definition}[Transaction]
\index{transaction}
Une \textbf{transaction} est une s\'equence d'op\'erations de lecture et
d'\'ecriture sur la base de donn\'ees, d\'elimit\'ee par un \texttt{BEGIN}
et termin\'ee par un \texttt{COMMIT} (validation) ou un \texttt{ROLLBACK}
(annulation).
\end{definition}

\begin{minted}{python}
# Exemple en Python avec psycopg2
import psycopg2

conn = psycopg2.connect("dbname=banque")
cur = conn.cursor()
try:
    cur.execute("UPDATE comptes SET solde = solde - 100 WHERE id = 1")
    cur.execute("UPDATE comptes SET solde = solde + 100 WHERE id = 2")
    conn.commit()   # Les deux operations sont validees ensemble
except Exception:
    conn.rollback()  # Aucune modification n'est appliquee
finally:
    cur.close()
    conn.close()
\end{minted}

% ------------------------------------------------------------
\section{Propri\'et\'es ACID}
% ------------------------------------------------------------

\begin{definition}[Propri\'et\'es ACID]
\index{ACID}
Les quatre propri\'et\'es fondamentales d'une transaction sont~:
\begin{enumerate}
  \item \textbf{Atomicit\'e} --- Toutes les op\'erations sont ex\'ecut\'ees ou aucune.
  \item \textbf{Coh\'erence} --- La transaction am\`ene la base d'un \'etat
        coh\'erent \`a un autre \'etat coh\'erent.
  \item \textbf{Isolation} --- Les transactions concurrentes ne voient pas
        les modifications interm\'ediaires des autres.
  \item \textbf{Durabilit\'e} --- Une fois valid\'ee (\texttt{COMMIT}), les
        modifications survivent \`a toute panne.
\end{enumerate}
\end{definition}

\begin{remarque}
L'atomicit\'e et la durabilit\'e sont assur\'ees par le journal de
transactions (\emph{Write-Ahead Log}, WAL). La coh\'erence est assur\'ee
par les contraintes d'int\'egrit\'e. L'isolation est assur\'ee par le
contr\^ole de concurrence.
\end{remarque}

% ------------------------------------------------------------
\section{Probl\`emes li\'es \`a la concurrence}
% ------------------------------------------------------------

Sans contr\^ole, l'ex\'ecution concurrente peut produire des anomalies~:

\begin{definition}[Anomalies de concurrence]
\index{anomalies de concurrence}
\begin{itemize}
  \item \textbf{Lecture sale} (\emph{dirty read}) --- $T_2$ lit une donn\'ee
        modifi\'ee par $T_1$ qui n'a pas encore valid\'e.
  \item \textbf{Lecture non r\'ep\'etable} (\emph{non-repeatable read}) ---
        $T_1$ lit une donn\'ee deux fois et obtient des valeurs diff\'erentes
        car $T_2$ l'a modifi\'ee entre-temps.
  \item \textbf{Lecture fant\^ome} (\emph{phantom read}) --- $T_1$ ex\'ecute
        une requ\^ete deux fois et obtient des lignes diff\'erentes car $T_2$
        a ins\'er\'e ou supprim\'e des lignes.
  \item \textbf{\'Ecriture perdue} (\emph{lost update}) --- Deux transactions
        lisent la m\^eme valeur, la modifient et l'\'ecrivent~: une des
        modifications est perdue.
\end{itemize}
\end{definition}

% ------------------------------------------------------------
\section{Niveaux d'isolation}
% ------------------------------------------------------------

\begin{definition}[Niveaux d'isolation SQL]
\index{niveaux d'isolation}
Le standard SQL d\'efinit quatre niveaux d'isolation~:
\begin{center}
\begin{tabular}{lccc}
\hline
\textbf{Niveau} & \textbf{Dirty read} & \textbf{Non-repeatable} & \textbf{Phantom} \\
\hline
\texttt{READ UNCOMMITTED} & Possible & Possible & Possible \\
\texttt{READ COMMITTED}   & Non      & Possible & Possible \\
\texttt{REPEATABLE READ}  & Non      & Non      & Possible \\
\texttt{SERIALIZABLE}     & Non      & Non      & Non \\
\hline
\end{tabular}
\end{center}
\end{definition}

\begin{attention}
Le niveau par d\'efaut varie selon le SGBD~: \texttt{READ COMMITTED} pour
PostgreSQL et Oracle, \texttt{REPEATABLE READ} pour MySQL/InnoDB. Choisir
un niveau trop faible expose aux anomalies~; un niveau trop fort r\'eduit
les performances.
\end{attention}

\begin{minted}{python}
-- Fixer le niveau d'isolation en SQL
SET TRANSACTION ISOLATION LEVEL SERIALIZABLE;
BEGIN;
    SELECT solde FROM comptes WHERE id = 1;
    UPDATE comptes SET solde = solde - 100 WHERE id = 1;
COMMIT;
\end{minted}

% ------------------------------------------------------------
\section{Th\'eorie de la s\'erialisabilit\'e}
% ------------------------------------------------------------

\begin{definition}[Ordonnancement (schedule)]
\index{ordonnancement}
Un \textbf{ordonnancement} $S$ est une s\'equence d'op\'erations
(lectures, \'ecritures, commits, rollbacks) provenant de plusieurs
transactions, qui pr\'eserve l'ordre interne de chaque transaction.
\end{definition}

\begin{definition}[S\'erialisabilit\'e par conflit]
\index{s\'erialisabilit\'e}
Deux op\'erations sont en \textbf{conflit} si elles portent sur le m\^eme
objet et au moins l'une est une \'ecriture. Un ordonnancement $S$ est
\textbf{s\'erialisable par conflit} s'il est \'equivalent par conflit \`a un
ordonnancement s\'eriel.
\end{definition}

\begin{theoreme}[Test de s\'erialisabilit\'e par graphe de pr\'ec\'edence]
\index{graphe de pr\'ec\'edence}
Construire le \textbf{graphe de pr\'ec\'edence} $G = (V, E)$ o\`u $V$ est
l'ensemble des transactions et $(T_i, T_j) \in E$ si une op\'eration de
$T_i$ est en conflit avec une op\'eration de $T_j$ et la pr\'ec\`ede.
L'ordonnancement est s\'erialisable par conflit si et seulement si $G$ est
acyclique.
\end{theoreme}

\begin{exemple}
Soit $S : r_1(A),\, w_1(A),\, r_2(A),\, w_2(A),\, r_1(B),\, w_1(B),\, r_2(B),\, w_2(B)$.
Les conflits sont~: $w_1(A) \to r_2(A)$, $w_1(A) \to w_2(A)$,
$w_1(B) \to r_2(B)$, $w_1(B) \to w_2(B)$. Le graphe a une seule ar\^ete
$T_1 \to T_2$, donc $S$ est s\'erialisable (\'equivalent \`a $T_1, T_2$).
\end{exemple}

% ------------------------------------------------------------
\section{Verrouillage \`a deux phases (2PL)}
% ------------------------------------------------------------

\begin{definition}[Protocole 2PL]
\index{verrouillage \`a deux phases}
Le protocole de \textbf{verrouillage \`a deux phases} (2PL) impose que
chaque transaction passe par deux phases~:
\begin{enumerate}
  \item \textbf{Phase de croissance}~: la transaction acquiert des verrous
        (partag\'es pour lecture, exclusifs pour \'ecriture) sans en lib\'erer.
  \item \textbf{Phase de d\'ecroissance}~: la transaction lib\`ere des verrous
        sans en acqu\'erir de nouveaux.
\end{enumerate}
\end{definition}

\begin{theoreme}[Correction de 2PL]
Tout ordonnancement produit par le protocole 2PL est s\'erialisable par conflit.
\end{theoreme}

\begin{definition}[2PL strict et rigoureux]
\begin{itemize}
  \item \textbf{2PL strict}~: les verrous exclusifs ne sont lib\'er\'es qu'au
        \texttt{COMMIT}/\texttt{ROLLBACK}. \'Evite les annulations en cascade.
  \item \textbf{2PL rigoureux}~: \emph{tous} les verrous sont lib\'er\'es au
        \texttt{COMMIT}/\texttt{ROLLBACK}. C'est le mode le plus courant.
\end{itemize}
\end{definition}

% ------------------------------------------------------------
\section{D\'etection et pr\'evention des interblocages}
% ------------------------------------------------------------

\begin{definition}[Interblocage (deadlock)]
\index{interblocage}
Un \textbf{interblocage} se produit lorsque deux ou plusieurs transactions
s'attendent mutuellement, chacune d\'etenant un verrou demand\'e par l'autre.
\end{definition}

\begin{proposition}[D\'etection par graphe d'attente]
Construire le \textbf{graphe d'attente}~: un arc $(T_i, T_j)$ signifie que
$T_i$ attend un verrou d\'etenu par $T_j$. Il y a interblocage si et
seulement si le graphe contient un cycle.
\end{proposition}

Strat\'egies de r\'esolution~:
\begin{itemize}
  \item \textbf{D\'etection p\'eriodique}~: v\'erifier les cycles et annuler
        la transaction la moins co\^uteuse (\emph{victim selection}).
  \item \textbf{Pr\'evention par estampilles}~: \emph{Wait-Die} ou
        \emph{Wound-Wait}.
  \item \textbf{Timeout}~: annuler toute transaction qui attend trop longtemps.
\end{itemize}

% ------------------------------------------------------------
\section{Contr\^ole de concurrence multi-versions (MVCC)}
% ------------------------------------------------------------

\begin{definition}[MVCC]
\index{MVCC}
Le \textbf{MVCC} (\emph{Multi-Version Concurrency Control}) maintient
plusieurs versions de chaque ligne. Une lecture ne bloque jamais une
\'ecriture et vice versa~:
\begin{itemize}
  \item Chaque transaction voit un \emph{snapshot} coh\'erent de la base
        au moment de son d\'ebut.
  \item Les \'ecritures cr\'eent de nouvelles versions sans \'ecraser les anciennes.
  \item Le ramasse-miettes (\emph{vacuum}) supprime les versions obsol\`etes.
\end{itemize}
\end{definition}

\begin{remarque}
PostgreSQL utilise MVCC pour tous les niveaux d'isolation. MySQL/InnoDB
utilise MVCC pour \texttt{READ COMMITTED} et \texttt{REPEATABLE READ},
et des verrous suppl\'ementaires pour \texttt{SERIALIZABLE}.
\end{remarque}

\begin{exemple}
Sous MVCC, la transaction $T_1$ peut lire la version $v_1$ d'une ligne
pendant que $T_2$ \'ecrit la version $v_2$~: aucun verrouillage n'est
n\'ecessaire pour la lecture.
\begin{minted}{python}
-- PostgreSQL : observer le MVCC
BEGIN;
SELECT xmin, xmax, * FROM comptes WHERE id = 1;
-- xmin = ID de la transaction qui a cree cette version
-- xmax = ID de la transaction qui l'a supprimee (0 si active)
COMMIT;
\end{minted}
\end{exemple}

% ------------------------------------------------------------
\section{Journal de transactions (WAL)}
% ------------------------------------------------------------

\begin{definition}[Write-Ahead Logging]
\index{WAL}
Le \textbf{WAL} (\emph{Write-Ahead Log}) est un journal s\'equentiel o\`u
toute modification est enregistr\'ee \emph{avant} d'\^etre appliqu\'ee aux
pages de donn\'ees. En cas de panne, le journal permet~:
\begin{itemize}
  \item \textbf{REDO}~: rejouer les modifications des transactions valid\'ees.
  \item \textbf{UNDO}~: annuler les modifications des transactions non valid\'ees.
\end{itemize}
\end{definition}

% ------------------------------------------------------------
\section{Formules cl\'es}
% ------------------------------------------------------------

\begin{formules}
\begin{itemize}
  \item \textbf{ACID}~: Atomicit\'e, Coh\'erence, Isolation, Durabilit\'e
  \item \textbf{S\'erialisabilit\'e}~: graphe de pr\'ec\'edence acyclique $\Leftrightarrow$ s\'erialisable
  \item \textbf{2PL}~: phase de croissance $\to$ point de verrouillage $\to$ phase de d\'ecroissance
  \item \textbf{Deadlock}~: cycle dans le graphe d'attente
  \item \textbf{MVCC}~: lectures non bloquantes via les snapshots multi-versions
\end{itemize}
\end{formules}

% ------------------------------------------------------------
\section{Exercices}
% ------------------------------------------------------------

\begin{exercice}
Soit l'ordonnancement $S : r_1(A),\, r_2(B),\, w_1(B),\, w_2(A)$.
Construire le graphe de pr\'ec\'edence. $S$ est-il s\'erialisable~?
\end{exercice}

\begin{exercice}
Donner un exemple d'interblocage entre deux transactions portant sur
deux comptes bancaires. Proposer une solution utilisant un ordre fixe
d'acquisition des verrous.
\end{exercice}

\begin{exercice}
\'Ecrire un script Python qui d\'emontre une lecture fant\^ome sous
\texttt{READ COMMITTED} dans PostgreSQL, puis montrer qu'elle dispara\^it
sous \texttt{SERIALIZABLE}.
\end{exercice}

\begin{exercice}
Expliquer pourquoi le 2PL strict \'evite les annulations en cascade.
Donner un contre-exemple avec le 2PL de base.
\end{exercice}

\begin{exercice}
Dans un syst\`eme MVCC, une transaction $T_1$ lit la ligne $r$ (version $v_1$),
puis $T_2$ modifie $r$ et valide (version $v_2$), puis $T_1$ relit $r$.
Quelle version voit $T_1$ sous \texttt{READ COMMITTED}~? Sous
\texttt{REPEATABLE READ}~? Justifier.
\end{exercice}
